% Options for packages loaded elsewhere
\PassOptionsToPackage{unicode}{hyperref}
\PassOptionsToPackage{hyphens}{url}
\documentclass[
]{article}
\usepackage{xcolor}
\usepackage{amsmath,amssymb}
\setcounter{secnumdepth}{-\maxdimen} % remove section numbering
\usepackage{iftex}
\ifPDFTeX
  \usepackage[T1]{fontenc}
  \usepackage[utf8]{inputenc}
  \usepackage{textcomp} % provide euro and other symbols
\else % if luatex or xetex
  \usepackage{unicode-math} % this also loads fontspec
  \defaultfontfeatures{Scale=MatchLowercase}
  \defaultfontfeatures[\rmfamily]{Ligatures=TeX,Scale=1}
\fi
\usepackage{lmodern}
\ifPDFTeX\else
  % xetex/luatex font selection
\fi
% Use upquote if available, for straight quotes in verbatim environments
\IfFileExists{upquote.sty}{\usepackage{upquote}}{}
\IfFileExists{microtype.sty}{% use microtype if available
  \usepackage[]{microtype}
  \UseMicrotypeSet[protrusion]{basicmath} % disable protrusion for tt fonts
}{}
\makeatletter
\@ifundefined{KOMAClassName}{% if non-KOMA class
  \IfFileExists{parskip.sty}{%
    \usepackage{parskip}
  }{% else
    \setlength{\parindent}{0pt}
    \setlength{\parskip}{6pt plus 2pt minus 1pt}}
}{% if KOMA class
  \KOMAoptions{parskip=half}}
\makeatother
\usepackage{longtable,booktabs,array}
\usepackage{calc} % for calculating minipage widths
% Correct order of tables after \paragraph or \subparagraph
\usepackage{etoolbox}
\makeatletter
\patchcmd\longtable{\par}{\if@noskipsec\mbox{}\fi\par}{}{}
\makeatother
% Allow footnotes in longtable head/foot
\IfFileExists{footnotehyper.sty}{\usepackage{footnotehyper}}{\usepackage{footnote}}
\makesavenoteenv{longtable}
\setlength{\emergencystretch}{3em} % prevent overfull lines
\providecommand{\tightlist}{%
  \setlength{\itemsep}{0pt}\setlength{\parskip}{0pt}}
\usepackage{bookmark}
\IfFileExists{xurl.sty}{\usepackage{xurl}}{} % add URL line breaks if available
\urlstyle{same}
\hypersetup{
  hidelinks,
  pdfcreator={LaTeX via pandoc}}

\author{}
\date{}

\begin{document}

\section{Commute brief --- lit review + answers
(2026-07-16)}\label{commute-brief-lit-review-answers-2026-07-16}

Reading order: \textbf{§0 first} (the Thea papers change the competitive
picture and answer several of your questions with hard numbers). Then
the Q\&A sections. Items tagged \textbf{{[}LIT PENDING{]}} are being
checked by background research agents; I'll append their findings.

Saved papers: \texttt{spf\_prototype/docs/papers/thea\_2512.08027.pdf}
(Helios overview, Swanson et al.~Dec 2025) and
\texttt{thea\_helios\_leadlithium\_blanket.pdf} (Pasmann et al., the
neutronics paper). The Thea neutronics \emph{slideshow} is behind a
Cloudflare bot-wall; curl can't fetch it --- an agent is trying
WebFetch.

\begin{center}\rule{0.5\linewidth}{0.5pt}\end{center}

\subsection{§0. The Thea papers --- this is the most important thing to
read}\label{the-thea-papers-this-is-the-most-important-thing-to-read}

The Helios blanket paper (Pasmann, van Riel, Khera, Romano, Koen,
Donovan, Swanson, \textbf{Gates}; and the overview has \textbf{E.J.
Paul, D.A. Gates, S. Pasmann}) is essentially a parallel effort in
\emph{exactly} our space, using the \emph{same stack} (OpenMC + DAGMC +
weight windows). Read it as both a validation and a competitive clock.

\textbf{What they did (so we DON'T claim it):} - A \textbf{stellarator
neutron source in OpenMC} (§3.2): a point cloud of discrete
\texttt{IndependentSource} sites, \textbf{isotropic} 14.1 MeV, with
\textbf{realistic Bosch--Hale fusion-reactivity weighting}
\texttt{s\_i\ =\ I\_i\ V\_i\ /\ (I\_T\ V\_T)} (core-peaked, Fig. 5). →
``native stellarator source in OpenMC'' is \textbf{not} novel. Our
differentiators: \textbf{spin polarization} (angular emission) and a
\textbf{compiled continuous sampler} vs.~their discrete cloud. -
\textbf{3D DAGMC forward neutronics} with \textbf{global variance
reduction}: they use \textbf{The Random Ray Method (TRRM) + FW-CADIS}
for neutron weight windows (§3.3), noting it beats MAGIC and gives
uniform uncertainty in deep-shield low-flux regions. \textbf{Photon
transport is not yet in TRRM} → photon WW via MAGIC. - Forward
\textbf{NWL map} (Fig. 10, poloidal×toroidal, Boozer coords): mean 0.95
MW/m², \textbf{peak 1.86 → peaking ≈ 2×}; ``comparable to ARIES-CS NWL
peaking \textasciitilde2× and dynamic factor \textasciitilde10×.'' -
\textbf{First-wall DPA} with their \textbf{cell-under-voxel (CUV)}
method (\texttt{MeshMaterialFilter}): mean 11.7 DPA/FPY, min/max
4.6/19.0 → FW life 17.7 FPY (10.5--42.2). Non-uniform (higher on broad
faces). Their proposed fix is \textbf{sector replacement}, not source
reshaping. - \textbf{Coil nuclear heating} (§3.6) --- \emph{this is our
problem, published and unsolved}: - Design limit \textbf{150 W/m³}
(cooling power). Binding constraint is \textbf{HEATING, not fluence.} -
Mean coil heating \textbf{80.6 W/m³}, \textbf{max 179.3 ± 4.2 W/m³ =
2.2× the mean}, \textbf{8 of 324 coils exceed}, concentrated
\textbf{inboard}. Stated remedy: ``additional cooling or
\textbf{shielding} considerations.'' - (Fluence limit is separately met:
min shaping-coil life 40.3 FPY.)

\textbf{What they explicitly DEFER (= our contribution):} - Blanket is a
\textbf{uniform radial offset} (§3.1, ``incrementally and uniformly
offset''). - §2.3: ``\textbf{Future work will include additional
shielding optimization\ldots{} non-uniform radial thicknesses.}''

\textbf{Their radial build (Table 2)} --- near-identical philosophy to
our corrected build, thinner total (\textasciitilde100 cm): FW 2.05 cm
(V-4Cr-4Ti) · breeder \textbf{50 cm Pb-17Li} (86.2\%
PbLi/EUROFER/He/SiC), Li6 \textbf{65\%} · interior shield \textbf{23 cm
WC + 7 cm B4C} · plasma vessel \textasciitilde13 cm (SS316L + vac +
borated water) · exterior \textbf{5 cm borated poly}. TBR target
\textbf{1.3}, met at Li6 65--70\% (3D: 1.307). Power multiplication 1.2.
Note: their device is \textbf{QA, R=8 m, A=4.5, B0=6 T, Bmax-coil 20 T,
min plasma--coil 1.2 m}, PbLi not FLiBe, 12 planar encircling + 324
shaping coils.

\textbf{Bottom line for framing:} Thea (Gates/Paul/Pasmann) is
sophisticated, active, and publishing in our exact niche as of Dec 2025.
They have a \textbf{real, concentrated, inboard coil-heating hotspot
(2.2× peaking, 8/324 over limit)} and \textbf{no spatial shield
solution} --- they defer it. Our closed-loop non-uniform shield
optimizer, \textbf{driven by coil hotspots}, is precisely the deferred
piece. But we must be \textbf{surgical about novelty}: claim
polarization, attribution/adjoint, and the closed spatial loop --- NOT
``stellarator source in OpenMC'' or ``3D stellarator neutronics,'' which
they (and ParaStell) already have.

\begin{center}\rule{0.5\linewidth}{0.5pt}\end{center}

\subsection{§1. Why is neutronics almost always downstream? What do
people actually
do?}\label{why-is-neutronics-almost-always-downstream-what-do-people-actually-do}

Short answer, and the Thea papers illustrate it perfectly: \textbf{the
design pipeline is plasma → coils → (freeze) → neutronics evaluation →
hand-iterate}, and neutronics is kept out of the plasma/coil optimizer
for a few compounding reasons:

\begin{enumerate}
\def\labelenumi{\arabic{enumi}.}
\tightlist
\item
  \textbf{Cost/differentiability mismatch.} Plasma+coil optimization
  (SIMSOPT/DESC/STELLOPT) runs thousands of cheap, differentiable
  evaluations of physics figures of merit (quasisymmetry, MHD,
  transport). A single 3D Monte Carlo neutronics evaluation with
  deep-shield weight windows is minutes-to-hours and \textbf{non-
  differentiable} (stochastic). You cannot put it inside a
  gradient-based shape loop without a surrogate.
\item
  \textbf{Weak \emph{shape} coupling, strong \emph{scenario} coupling.}
  To leading order the neutron \textbf{source} is set by the fusion rate
  (power, profiles), and the wall/coil load is set by geometry+standoff
  --- both of which are only loosely sensitive to the fine plasma-shape
  knobs the optimizer is actually turning. So neutronics is treated as
  separable and checked afterward.
\item
  \textbf{Tooling separation.} Different codes, teams, meshes.
  ParaStell/Stellarmesh/OpenMC live downstream of DESC/SIMSOPT by
  construction.
\end{enumerate}

\textbf{What they do instead of co-optimizing:} they \textbf{design to a
limit and add bulk shielding / margin.} Two strategies, both visible in
the Thea papers: - \textbf{Buy margin with standoff} (Helios
philosophy): make the plasma \emph{less strongly shaped} so coils sit
further out (1.2 m), leaving room for a \textbf{uniform} thick
blanket+shield. Overview paper, verbatim: \emph{``The Helios equilibrium
is less strongly shaped, allowing the coils to be further away. This
significantly eases the design of the breeding blanket while protecting
the coils\ldots{} ARIES-CS required a highly optimized non-uniform
blanket that prioritizes shielding at the expense of breeding zone in
certain places.''} - \textbf{Spatially optimize the blanket by hand}
(ARIES-CS philosophy): strongly-shaped/compact → not enough uniform room
→ thin the breeder and insert shield at the hot spots, \textbf{by
engineering judgment}.

So yes --- essentially \textbf{``add enough shielding to meet the peak
limit, iterate the plasma by hand if you can't.''} Peaking factor
\textbf{is} effectively the binding constraint (it sets how thick the
shield must be to protect the \emph{worst} coil), but nobody closes it
in an automated spatial loop. That gap --- \textbf{automated,
spatially-resolved, hotspot-driven} shield allocation --- is the
opening, and it exists precisely \emph{because} MC-in-the-loop is
expensive and non-differentiable (which is why our cheap surrogate +
adjoint engine matters). \textbf{{[}LIT PENDING: confirming nobody
co-optimizes neutronics peaking with plasma/coil shape --- agent
running.{]}}

\begin{center}\rule{0.5\linewidth}{0.5pt}\end{center}

\subsection{§2. Wall vs.~coil shielding; why sacrifice breeder; why
hasn't anyone done it for
stellarators}\label{wall-vs.-coil-shielding-why-sacrifice-breeder-why-hasnt-anyone-done-it-for-stellarators}

\textbf{Why thin the \emph{breeder} to make room for \emph{shield}
(rather than move coils or grow the machine)?} - \textbf{Machine cost
scales steeply with size.} Major radius drives everything (magnet stored
energy, structure, building, cost). Pushing coils out to make room is
the most expensive possible fix. The whole point of a \emph{compact}
stellarator reactor is to NOT do that. So the radial build between
plasma and coil is a \textbf{fixed, scarce budget}, and within it
breeder and shield compete. (This is the ARIES-CS vs Helios split
above.) - \textbf{Coils must stay close} because the field they must
produce at the plasma scales with distance --- Thea (blanket §2):
\emph{``The further away the shaping coils, the higher the magnetic
field required on the coil, which incentivizes a minimum distance
between coils and plasma.''} Move coils out → need higher on-coil field
→ bigger/more-expensive HTS, higher stress. So standoff is bounded from
both sides. - \textbf{Coils are the irreplaceable, lifetime-limiting
component} (40 yr), whereas the \textbf{first wall is a consumable}
(Thea FW life \textasciitilde17.7 FPY, and they plan \textbf{sector
replacement}). So you spend breeder margin to protect the thing you
can't replace. - \textbf{Maintenance drives architecture} (Helios:
sector removal between planar coils, 84-day biennial outage, 88\%
capacity factor) --- but that's about \emph{access}, not about the
breeder/shield trade per se.

\textbf{Why ``wall shielding'' isn't a thing but ``coil shielding'' is}
(your instinct, confirmed): you can't put shield \emph{in front of} the
first wall --- it \emph{is} the first surface. First-wall load is
reduced only by reshaping the \textbf{source} (plasma shape,
polarization) or the \textbf{wall} (shape/armor) or by
\textbf{replacing} it. The coil, sitting behind \textasciitilde1 m of
blanket, is the only place a \textbf{shield} lever exists. Hence:
\textbf{first wall = source/ geometry problem; magnets = shield
problem.}

\textbf{Why hasn't anyone automated it for stellarators?} Three reasons,
all now evidenced: 1. \textbf{It's genuinely hard and expensive}
(MC-in-loop, non-differentiable, 3D CAD geometry regen per iterate). 2.
\textbf{The tooling only just matured} --- ParaStell (2024) and
Stellarmesh (Thea) made parametric 3D CAD→DAGMC feasible; before that
you couldn't cheaply regenerate a non-uniform stellarator blanket. 3.
\textbf{The people who could just got here} --- ParaStell
\textbf{defers} the loop; Thea \textbf{defers} non-uniform thickness.
Two independent groups parked at the same ``future work.'' That's the
strongest possible signal the automated spatial loop is unclaimed
\emph{and wanted}. \textbf{{[}LIT PENDING: agent confirming no closed
loop.{]}}

\begin{center}\rule{0.5\linewidth}{0.5pt}\end{center}

\subsection{§3. A principled, plasma-power-specific attribution cutoff
(not
5\%)}\label{a-principled-plasma-power-specific-attribution-cutoff-not-5}

Your instinct is right and the Thea numbers make it concrete. The cutoff
should be the \textbf{engineering limit}, which \emph{is}
plasma-power-specific because the loads scale with fusion power:

\begin{itemize}
\tightlist
\item
  \textbf{The load is power-normalized.} NWL and coil heating scale
  \textasciitilde linearly with fusion power (W/m² and W/m³ per MW of
  fusion). Thea normalizes everything ``per 958 MW.'' So a device at 2×
  the power density has \textasciitilde2× the loads for the same
  geometry → the \emph{same geometry} can be under-limit at low power
  and over-limit at high power. The attribution cutoff = ``\textbf{where
  does load exceed the fixed material limit at THIS device's power}.''
  That's inherently power- and device-specific, not a fixed percentile.
\item
  \textbf{Concretely (Thea):} limit 150 W/m³; at 958 MW, \textbf{8/324
  coils} exceed. At higher power more coils cross; at lower power none
  do. \textbf{The over-limit SET is the attribution target}, and its
  size is an output, not an assumption. This is strictly better than
  ``top 5\%.''
\end{itemize}

\textbf{``Do they just build thicker walls, and how does that affect
magnet distance?''} --- Exactly the tension: thicker shield → either (a)
push coils out (expensive, higher on-coil field) or (b) thin the breeder
(lose TBR) at fixed envelope. Thea buys margin by \textbf{standoff} (1.2
m, uniform build); ARIES-CS buys it by \textbf{spatial breeder→shield
trade}. Our optimizer automates (b) \emph{only where needed} (over the 8
hot coils), spending the \emph{minimum} breeder --- which is the
marginal-value idea below.

\textbf{Marginal-value weighting (you liked this):} weight each
over-limit location by how \emph{binding} it is ---
\texttt{d(lifetime)/d(load)} or how much relaxing it moves the limiting
constraint. For a \textbf{sharp-threshold} failure (insulator dose,
quench) the near-threshold coils dominate; for a \textbf{wear-out} mode
(fluence accumulation) the integral matters. Thea's coil case is a
\textbf{hard cooling-power threshold} (150 W/m³) → near-threshold/over-
threshold coils dominate → attribute and shield exactly those 8. This
turns ``top 5\%'' into ``the set whose marginal protection is worth the
breeder it costs.'' \textbf{{[}LIT PENDING: whether
marginal-value/limit-based attribution cutoff is used anywhere in fusion
neutronics.{]}}

\begin{center}\rule{0.5\linewidth}{0.5pt}\end{center}

\subsection{§4. Device-to-device (QA/QH/QI) variation + concentration
index}\label{device-to-device-qaqhqi-variation-concentration-index}

You want to (a) map how attribution region size / concentration varies
across the QS zoo, and (b) formalize a Gini/participation concentration
index, with the thesis \emph{``attribution concentration predicts
shield-optimizer efficacy.''} I think this is a real, under-explored
science contribution and it connects cleanly to our saturation result:

\begin{itemize}
\tightlist
\item
  \textbf{Why devices differ:} QA (tokamak-like, weaker shaping, Helios)
  vs QH (stronger helical shaping, our device) vs QI have very different
  \textbar B\textbar{} spectra, standoff variation, and boundary
  curvature → different load distributions. Thea's QA coil hotspot is
  \textbf{inboard-concentrated} (8/324). Our QH baseline was a
  \textbf{localized consecutive cluster} (coils 27--30). Different
  concentration signatures already, n=2.
\item
  \textbf{The concentration index} (Gini / participation ratio of the
  culprit field \texttt{C\_v}, or of the coil-load distribution) makes
  ``concentrated vs diffuse'' a single number. Then the thesis is
  testable across a QUASR batch: does high concentration ⇒ localized
  shield works well (few patches, low breeder cost), and low
  concentration ⇒ localized shield saturates fast / you need broad
  coverage or a source-side fix?
\item
  \textbf{Ties to our result:} we \emph{measured} that a localized patch
  saturates on a leak-around floor. Concentration should predict
  \emph{where that floor sits} and \emph{how many patches} you need ---
  the quantitative bridge between attribution (diagnosis) and shield
  optimization (cure). \textbf{{[}LIT PENDING: is QS-zoo
  load-distribution variation studied? Any Gini/participation
  concentration metric in fusion neutronics? --- agent running.{]}}
\end{itemize}

Note on your aside ``\emph{higher-concentration devices may be the
unpolarized ones}'': worth testing directly --- polarization
redistributes the source, which can either sharpen or spread the coil
load depending on B geometry. Don't assume; it's a measurement (and a
nice result either way).

\begin{center}\rule{0.5\linewidth}{0.5pt}\end{center}

\subsection{§5. Slant path + your scattering question (this is a sharp
question and it refines the
idea)}\label{slant-path-your-scattering-question-this-is-a-sharp-question-and-it-refines-the-idea}

You asked: does scattering randomize the neutron path, and what does
that do to the slant-path effective-length (the \texttt{t/cosθ} idea)?
Yes --- and thinking it through actually \emph{connects the slant-path
lever to our saturation result}:

\begin{itemize}
\tightlist
\item
  \textbf{The \texttt{t/cosθ} slant-path enhancement is an UNCOLLIDED
  (first-flight) effect.} It's the extra attenuation a neutron gets by
  entering the shield obliquely and traversing more material
  \emph{before its first collision}. It applies cleanly to the
  \textbf{direct/uncollided} component \texttt{exp(−Σ·t/cosθ)}.
\item
  \textbf{Once a neutron scatters, it forgets its entry angle.} The
  direction randomizes; the ``I entered at a slant'' memory is largely
  lost after \textasciitilde1 collision. So the scattered component's
  transmission is governed by the \textbf{removal cross section +
  buildup}, not by the original incidence angle.
\item
  \textbf{Scattering also \emph{lengthens} the effective path} via
  random-walk (total track length ≫ straight-line thickness) --- this is
  exactly why the \emph{removal} MFP is what matters and why buildup
  factors exist.
\item
  \textbf{The punchline, tied to our data:} our δ-sweep showed the coil
  flux past \textasciitilde1 mfp is \textbf{dominated by scattered /
  leak-around paths} (that's the saturation floor). Those paths have
  \emph{forgotten} the entry angle. So \textbf{the slant-path (channel
  2) lever of polarization mostly affects the uncollided/first-flight
  component, which is exactly the component that's already small deep
  behind thick shield.} Therefore: \textbf{polarization's coil-dose
  lever is probably dominated by channel 1 (source redistribution), with
  channel 2 (slant path) a smaller, shallow-shield / near-surface
  effect.} That's a non-obvious, honest refinement --- and it's
  \emph{testable}: decompose the coil-dose change from polarization into
  uncollided (angle-sensitive) vs scattered (angle-washed) components in
  the MC. If channel 2 is real it lives in the uncollided tally.
\end{itemize}

This also means the \textbf{first-wall} slant-path story is cleaner than
the coil one: the first wall sees mostly uncollided source neutrons
(transparent plasma), so incidence-angle effects survive there; deep at
the coil, scattering has washed them out.

\begin{center}\rule{0.5\linewidth}{0.5pt}\end{center}

\subsection{§6. Random Ray, the angular idea, and the adjoint
plan}\label{random-ray-the-angular-idea-and-the-adjoint-plan}

\textbf{Partial answer already from Thea's paper:} Random Ray
\textbf{is} in OpenMC, \textbf{is} coupled to \textbf{FW-CADIS} (an
adjoint-importance method), and Thea \emph{uses} it (TRRM) for weight
windows --- so the adjoint importance is \emph{already being computed}
in this stack, but for \textbf{variance reduction}, not \textbf{source
attribution}. Two gaps Thea names that map onto your PI's-student idea:
- \textbf{CORRECTION (agent found this --- do not overclaim):}
\emph{anisotropic / angular random ray is ALREADY PUBLISHED} ---
\textbf{Neame \& Cosgrove, ``Linear Sources and Anisotropic Scattering
in the Random Ray Method,'' NSE 2024, DOI 10.1080/00295639.2024.2394729}
--- flat→linear source and isotropic→\textbf{anisotropic scattering up
to P3}, implemented in \textbf{SCONE} (not OpenMC). So ``anisotropic
random ray'' as a method is \textbf{not novel.} What \emph{is} still
open, and is the defensible engine claim: 1. It does \textbf{not exist
in OpenMC's} random ray (OpenMC RR is MG-only, \textbf{isotropic source,
scalar/angle- integrated flux, no anisotropic scattering}) --- porting
it in is engineering-novel; 2. retaining an \textbf{angularly-resolved
adjoint importance field} \texttt{ψ†(r,Ω,E)} as the \emph{solution
output} (not merely an anisotropic scattering source that still
collapses to scalar flux) is under-explored even in SCONE, and is
exactly what a \textbf{streaming-sensitive coil} attribution needs
(precedent: FW/CADIS-Ω, arXiv:1612.00793, shows adjoint \emph{angular}
flux matters for deep-penetration/streaming); 3. the
\textbf{application} --- angular adjoint importance → stellarator
\textbf{coil source-attribution} --- is unclaimed. So the methods-paper
engine claim is (ii)+(iii), NOT ``we invented anisotropic random ray.''
- \textbf{OpenMC has no general continuous-energy adjoint solver}; the
accepted routes are (a) \textbf{multigroup adjoint} (what
random-ray/FW-CADIS effectively do) and (b) \textbf{reciprocity}:
forward-transport a pseudo- source placed \emph{at the coil} into the
shield and tally where it reaches the plasma = the
contributon/importance map. Both are implementable in-stack.

\textbf{Your ``do both to validate'' plan is exactly right} and is
standard practice: forward MC with \textbf{source-region tagging} on one
reference config = ground truth; adjoint/contributon importance = the
cheap engine; show they agree; then run the adjoint for the sweep.
\textbf{{[}LIT PENDING: full Random Ray doc read +
adjoint-in-stellarator novelty + whether angular random ray is being
pursued --- agent running; I'll append the doc-level details.{]}}

\textbf{{[}LIT PENDING: ``first adjoint/contributon coil-load
source-attribution in a 3D stellarator'' novelty check.{]}}

\begin{center}\rule{0.5\linewidth}{0.5pt}\end{center}

\subsection{§7. Real emissivity --- visualization, pluggable infra, and
the
factorization}\label{real-emissivity-visualization-pluggable-infra-and-the-factorization}

\begin{itemize}
\tightlist
\item
  \textbf{Thea already uses real emissivity} (Bosch--Hale, core-peaked,
  Fig. 5) for their \emph{forward} NWL. So our uniform-emissivity
  culprit map is a \emph{simplification relative to Thea's forward
  source}. That's fine --- but it means we should (a) build the
  pluggable-emissivity infra and (b) USE realistic profiles for the
  \emph{science}, keeping uniform only for the clean decomposition. This
  makes our tool a superset of theirs (adds attribution

  \begin{itemize}
  \tightlist
  \item
    polarization + pluggable scenario) rather than a subset.
  \end{itemize}
\item
  \textbf{Pluggable weight function} \texttt{r(ρ,θ,ζ)}: keep the
  geometry+polarization kernel fixed; users drop in a profile (or a
  kinetic-profile file). Clean software boundary; lets others study
  their scenarios.
\item
  \textbf{The 3D translucent-voxel visualization you sketched:} yes ---
  render the plasma as a volume of voxels whose \textbf{opacity ∝
  emissivity·contribution} and \textbf{color ∝
  contribution-to-peak-coil-load}, so low-culprit plasma is transparent
  and culprit regions ``light up'' as solid colored blobs inside the
  boundary; add \textbf{rotating views + a couple of poloidal/toroidal
  cross-section cut-planes} so you can see interior structure (core vs
  edge culprits). This is a genuinely better figure than the surface
  paint, and it's the natural way to show the \emph{volumetric} culprit
  field with realistic (core-weighted) emissivity. Cheap to build
  (matplotlib 3D scatter with alpha, or a proper volume render).
\item
  \textbf{The factorization}
  \texttt{C\_v\ ≈\ {[}geometry\ kernel{]}\ ×\ {[}polarization\ g{]}\ ×\ {[}emissivity\ r{]}}
  is the reason the analytic tool is the decomposition workhorse: the
  three factors enter \textbf{multiplicatively} in the free- streaming
  first-wall limit, so a \textbf{factorial study \{uniform, realistic\}
  × \{unpol, A, B/C\} × \{device zoo\}} cleanly separates how much of
  the culprit \emph{location} and \emph{concentration} is set by
  geometry vs shifted by polarization vs shifted by profile. Full
  transport breaks strict separability (scattering mixes terms), so:
  \textbf{analytic factorization = the clean decomposition; MC =
  validation + the deep-coil (scattered) regime.} Best-case finding: a
  stable trend like ``geometry sets location; polarization moves it by
  X\%; profile changes concentration not location'' --- a transferable
  result, not just a method.
\end{itemize}

\begin{center}\rule{0.5\linewidth}{0.5pt}\end{center}

\subsection{§8. Two-paper split --- I strongly agree, and the timing
makes it
urgent}\label{two-paper-split-i-strongly-agree-and-the-timing-makes-it-urgent}

Your instinct to split is right, and the Thea Dec-2025 paper makes the
\textbf{methods paper time-critical} (stake the tooling claims before
the active groups --- Thea/CNERG --- extend into them).

\begin{itemize}
\tightlist
\item
  \textbf{Paper 1 (Methods / tools):}

  \begin{enumerate}
  \def\labelenumi{\arabic{enumi}.}
  \tightlist
  \item
    compiled \textbf{polarized} \texttt{StellaratorSource} (+ SPF in
    \texttt{TokamakSource}) --- differentiate vs Thea's isotropic
    point-cloud and vs Bae's polarized tokamak source;
  \item
    \textbf{automatic non-uniform deformation} of an
    initially-constant-offset blanket (the thing ParaStell \& Thea both
    defer);
  \item
    the \textbf{adjoint/contributon attribution engine} (+ angular
    random-ray if your PI's student's idea lands), validated against
    forward source-tagged MC;
  \item
    the \textbf{analytic free-streaming factorization} tool + pluggable
    emissivity.
  \end{enumerate}
\item
  \textbf{Paper 2 (Science / across the QS zoo):} your 5-part spine ---
  attribution (culprit maps, limit-based cutoff, concentration index) →
  concentration-as-predictor → interventions (source-side
  polarization+shape for the wall; coverage-not-thickness shield for the
  coils, with the leak-around floor) → geometry/ polarization/scenario
  decomposition → adjoint-validated engine.
\end{itemize}

Honest note: Paper 1 must be \textbf{surgical on novelty} given
Thea/ParaStell. The claims that survive are polarization,
attribution/adjoint, closed-loop non-uniform shield, and (maybe) angular
random-ray --- NOT ``stellarator neutronics in OpenMC.''

\begin{center}\rule{0.5\linewidth}{0.5pt}\end{center}

\subsection{§9. Literature-search results (3 agents, cited) --- the
novelty
verdicts}\label{literature-search-results-3-agents-cited-the-novelty-verdicts}

Tags: \textbf{{[}EST{]}} read from primary/authoritative source ·
\textbf{{[}INF{]}} inferred from adjacent sources · \textbf{{[}NULL{]}}
searched hard, apparently unpublished · \textbf{{[}CNV{]}}
could-not-verify (paywalled/under-indexed). Numbers from abstract-level
snippets are flagged --- spot-check primaries before quoting in a paper.

\subsubsection{9.1 Why neutronics is downstream / co-optimization (agent
A)}\label{why-neutronics-is-downstream-co-optimization-agent-a}

\begin{itemize}
\tightlist
\item
  \textbf{{[}EST{]}} The stated reason is overwhelmingly \textbf{MC is
  \textasciitilde5 orders of magnitude too slow for an optimization
  loop} + it needs a \textbf{fully-built CAD device}. Lion et al.~2022
  (Nucl. Fusion 62 076040, DOI 10.1088/1741-4326/ac6a67) say resolving
  neutronics needs 1e8--1e9 histories, ``tractable only by HPC,'' and
  their deterministic surrogate is ``\textasciitilde5 orders faster.''
  Bogaarts \& Warmer 2024 (arXiv:2411.16369) echo it.
\item
  \textbf{{[}EST{]}} Practice is \textbf{limit-driven radial-build
  sizing}: set a peak-NWL / magnet-dose limit, then tailor (often
  non-uniform) blanket+shield to meet it. Peaking managed by
  \textbf{first-wall shaping} (Lion: peaking \textbf{1.69→1.23}) or
  \textbf{non-uniform shield} (ARIES-CS). STELLOPT/SIMSOPT/DESC
  objectives contain \textbf{no neutronics}.
\item
  \textbf{{[}EST/strong NULL{]}} Neutronics enters the \emph{inner}
  plasma/coil optimization loop \textbf{only via a geometric standoff
  proxy} --- min plasma--coil distance, or the
  \textbf{magnetic-gradient-scale-length} differentiable proxy (DOI
  10.1088/1741-4326/ae69fa). \textbf{NWL-peaking / magnet-dose as a
  transport-derived objective on plasma boundary or coil shape is NOT
  demonstrated} and is \textbf{repeatedly named as future work}
  (ParaStell: ``coupling ParaStell to machine-driven optimization is
  planned''; Lion: ``can then\ldots{} be used within an optimization
  framework''). → confirms §1/§2 above.
\end{itemize}

\subsubsection{9.2 Adjoint / contributon attribution + Random Ray (agent
B) --- the strongest
novelty}\label{adjoint-contributon-attribution-random-ray-agent-b-the-strongest-novelty}

\begin{itemize}
\tightlist
\item
  \textbf{{[}NULL{]}} \textbf{Adjoint/contributon backward attribution
  of first-wall OR coil load to plasma SOURCE REGIONS in a 3-D
  stellarator = unpublished.} Searches over CNERG/ParaStell, W7-X,
  HELIAS, Thea, ORNL/Denovo/Shift, contributon theory returned only (i)
  \emph{forward} NWL maps (Lion 2022, ParaStell), (ii) adjoint for
  \emph{tokamak diagnostics} (reciprocity, axisymmetric ---
  arXiv:1504.03073), (iii) \emph{geometry/MHD} adjoints (Paul/Antonsen).
  \textbf{Defensible novelty} if framed as \emph{backward source-region
  attribution} (contributon map) vs the abundant forward maps. Formal
  basis: contributon theory, Williams UWFDM-1338. Hedge with ``to our
  knowledge'' (proceedings under-indexed).
\item
  \textbf{{[}EST{]}} \textbf{OpenMC has NO continuous-energy adjoint.}
  Only \textbf{multigroup Random-Ray adjoint} (transpose scattering
  matrix + swap νΣ\_f↔χ), used for FW-CADIS weight windows. Workarounds:
  MG RR adjoint; \textbf{reciprocity} (detector-side forward source); CE
  perturbation/sensitivity (not an importance field); external S\_N
  (Denovo/ ADVANTG). CADIS/FW-CADIS = the mature hybrid VR standard
  (Wagner; MS-CADIS for ITER shutdown dose, CNERG).
\item
  \textbf{{[}EST{]}} \textbf{Random Ray = stochastic Method of
  Characteristics, multigroup only, isotropic source, angular dependence
  integrated out (scalar flux), no anisotropic scattering in OpenMC,
  flat/linear source; adjoint IS supported}
  (\texttt{settings.random\_ray{[}\textquotesingle{}adjoint\textquotesingle{}{]}=True}).
  Tramm et al.~NSE 2023.
\item
  \textbf{{[}EST{]} correction to §6:} anisotropic random ray is
  \textbf{already published} (Neame \& Cosgrove, NSE 2024, SCONE, P3).
  Open + defensible: angular-resolved adjoint \emph{output} field,
  porting anisotropy into OpenMC, and the stellarator-coil-attribution
  application.
\end{itemize}

\subsubsection{9.3 QS-zoo variation, concentration metrics, wall-vs-coil
trade (agent
C)}\label{qs-zoo-variation-concentration-metrics-wall-vs-coil-trade-agent-c}

\begin{itemize}
\tightlist
\item
  \textbf{{[}NULL{]}} \textbf{No systematic cross-configuration
  (QA/QH/QI) or QUASR-wide NWL-distribution/peaking study.} Lion 2022
  computes peaking per-config (HELIAS-3/4/5 + a QA case) but it's a
  \emph{method} paper, not a controlled zoo comparison. \textbf{QUASR
  has only QA+QH} (\textasciitilde370k configs; Giuliani,
  arXiv:2409.04826) --- a QI-inclusive comparison needs QI configs from
  elsewhere (W7-X line / Stellaris). Novelty = ``peaking computed
  device-by-device but never \textbf{mapped across the quasisymmetry
  landscape}.'' Do NOT claim to invent NWL peaking.
\item
  \textbf{{[}CNV→opportunity{]}} \textbf{Gini / Lorenz /
  participation-ratio concentration metrics are NOT used in fusion
  neutronics} --- only crude peak/mean. Importing them is \textbf{novel
  in application} (not in invention). Participation ratio (PR =
  (Σvᵢ²)²/Σvᵢ⁴ = ``effective number of hot spots'') is the most natural
  for ``how many patches carry the load.'' Frame as borrowed from
  economics (Gini) / condensed-matter (PR).
\item
  \textbf{{[}EST{]}} \textbf{Wall-vs-coil / breeder-for-shield trade is
  well documented.} ARIES-CS swapped blanket→WC shield at min-standoff
  spots, cutting local plasma-to-mid-coil build
  \textbf{\textasciitilde1.79 m → \textasciitilde1.31 m
  (\textasciitilde20--30\%)} while preserving \textbf{TBR ≈ 1.1}
  (El-Guebaly UWFDM-1336). Reason: machine size/cost scale steeply with
  standoff; coils must stay close (field-on-coil rises with distance).
  Corroborated by Helios (1.2 m budget) and ARC (FLiBe does triple
  duty).
\item
  \textbf{{[}NULL, but ``reviewer may call it obvious''{]}}
  \textbf{``Attribution concentration predicts localized-intervention
  effectiveness'' appears unpublished.} Defend it by (a) quantifying
  with a real concentration metric, (b) making it a \emph{falsifiable}
  cross-device prediction, (c) showing failure cases. \textbf{Honesty
  flag from the agent, which I second:} this is the same shape of claim
  as our earlier cheap-predictor that did NOT validate at n=5 (benefit
  tracked headroom, not shape) --- hold it to the same skeptical bar
  before asserting.
\end{itemize}

\subsubsection{9.4 Net novelty scorecard (what survives adversarial
review)}\label{net-novelty-scorecard-what-survives-adversarial-review}

\begin{longtable}[]{@{}
  >{\raggedright\arraybackslash}p{(\linewidth - 2\tabcolsep) * \real{0.5000}}
  >{\raggedright\arraybackslash}p{(\linewidth - 2\tabcolsep) * \real{0.5000}}@{}}
\toprule\noalign{}
\begin{minipage}[b]{\linewidth}\raggedright
Claim
\end{minipage} & \begin{minipage}[b]{\linewidth}\raggedright
Verdict
\end{minipage} \\
\midrule\noalign{}
\endhead
\bottomrule\noalign{}
\endlastfoot
First \textbf{native, reusable, general} stellarator \texttt{Source}
class in OpenMC (+ \textbf{polarization}) & \textbf{Novel} --- Thea has
\emph{a} stellarator source \emph{in} OpenMC but it is an internal
workflow emitting a \textbf{discrete point-cloud of IndependentSource
sites} (isotropic, scenario-specific), NOT a drop-in reusable Source
class; Bae's polarized source is tokamak. \textbf{Wording:} claim
``first native/reusable/general,'' NOT ``first stellarator source in
OpenMC.'' \textbf{To substantiate ``others can use,'' upstream it or
release a documented package.} \\
Automated \textbf{non-uniform} blanket/shield spatial optimization
(closed loop) & \textbf{Novel} (ParaStell + Thea both defer it) \\
\textbf{Adjoint/contributon coil-load → plasma-source attribution in 3D
stellarator} & \textbf{Novel {[}NULL{]}} --- strongest claim \\
Coverage-not-thickness + leak-around floor finding & \textbf{Novel} (our
measurement) \\
Concentration index (Gini/PR) for load/attribution + ``predicts
efficacy'' & \textbf{Novel in application}; the \emph{predictor} thesis
needs skeptical validation \\
Cross-QS-zoo peaking/concentration map & \textbf{Novel} (never mapped
across the landscape) \\
``Stellarator source in OpenMC'' / ``3D stellarator neutronics'' &
\textbf{NOT novel} (Thea, ParaStell) --- don't claim \\
``Anisotropic random ray'' as a method & \textbf{NOT novel} (Neame \&
Cosgrove 2024) --- don't claim \\
NWL peaking / magnet dose as plasma/coil optimization objective &
\textbf{{[}strong NULL{]}} --- future-work everywhere; a separate big
claim if you pursue it \\
\end{longtable}

\subsubsection{9.5 Still open}\label{still-open}

\begin{itemize}
\tightlist
\item
  Thea neutronics \textbf{slideshow} stayed Cloudflare-blocked (agent
  couldn't fetch it either). Not critical --- the two papers cover the
  substance. Grab it from your own browser if you want the slides.
\item
  Several ARIES-CS / Lion exact numbers are abstract-level; pull primary
  PDFs (Lion TU/e copy; El-Guebaly UWFDM-1336) before quoting figures in
  a manuscript.
\end{itemize}

\end{document}
